% =============================================================================
% 06_anhang.tex — Anhang
% =============================================================================
% Der Anhang enthält ergänzende Materialien, die den Lesefluss im Hauptteil
% gestört hätten. Typische Inhalte:
%   - Selbst erstellte Unterlagen (Schaltpläne, Datenblätter, Protokolle)
%   - Unveröffentlichte Quellen (interne Dokumente, E-Mails, Interviews)
%   - Detaillierte Auswertungen (Rohdaten, umfangreiche Tabellen)
%   - Vollständige Quelltexte (sofern nicht im Repo verlinkt)
% =============================================================================

% Anhang-Formatierung: Kapitel werden mit Buchstaben nummeriert (A, B, C …)
\appendix

\chapter{Anhang}
\label{chap:anhang}

% --- A. Schaltpläne und Hardware-Dokumentation --------------------------------
\section{Schaltpläne und Hardware-Dokumentation}
% Füge Schaltpläne, Stücklisten (BOM) oder Datenblätter ein.
% \begin{figure}[htbp]
%   \centering
%   \includegraphics[width=0.9\textwidth]{schaltplan_geotracker}
%   \caption{Schaltplan des Geotracker-Prototyps}
%   \label{fig:schaltplan}
% \end{figure}
<<Schaltpläne, Stücklisten (BOM) und weitere Hardware-Unterlagen hier einfügen.>>

% --- B. Messprotokolle und Testergebnisse ------------------------------------
\section{Messprotokolle und Testergebnisse}
% Detaillierte Tabellen mit Testfällen und Messwerten, die im Hauptteil
% zu umfangreich wären.
<<Vollständige Testprotokolle, Messwert-Tabellen und Auswertungen hier einfügen.>>

% --- C. Unveröffentlichte Quellen --------------------------------------------
\section{Unveröffentlichte Quellen}
% Interne Dokumente, E-Mails, Interview-Transkripte, die nicht öffentlich
% zugänglich sind, aber als Quellen dienen.
<<Unveröffentlichte Quellen (mit Datum und Kontext) hier einfügen.>>

% --- D. Weitere Unterlagen ---------------------------------------------------
\section{Weitere Unterlagen}
% Platzhalter für sonstige relevante Materialien.
<<Weitere ergänzende Unterlagen hier einfügen.>>
